\begin{problem}{N}{New Nickname}{1 second}

\noindent
Tired of her old nickname, Haibara has decided to invent a new one. Since she has a likeness for palindromes, she wants her new name to satisfy a set of constraints about the palindromic substrings.

\noindent
She gives you a binary string $P$ of length $N$. You must construct a string $S$ of length $N$ consisting of \textbf{lowercase or uppercase English letters} such that, for each integer $k$ from $1$ to $N$:

\begin{itemize}
\item $P_k = 1$ means that $S$ must contain a palindromic substring of length $k$.
\item $P_k = 0$ means that $S$ must not contain any palindromic substring of length $k$.
\end{itemize}

\noindent
A substring is a contiguous sequence of characters within the string, and a palindrome is a string that reads the same forwards and backwards. \textbf{Note that palindromes are case-sensitive; for example, "aA" and "Abba" are not considered palindromes.}

\noindent
Your task is to find any valid string $S$ that satisfies all the given requirements or determine that no such string exists.

\InputFile

\noindent
The first line of the input contains a single integer $T$ ($1 \le T \le 1000$) --- the number of testcases.

\noindent
Each testcase consists of two lines. The first line contains a single integer $N$ ($1 \le N \le 5 \cdot 10^5$), the length of the string.

\noindent
The second line contains a binary string $P$ of length $N$, representing the length restrictions for the palindromic substrings.

\noindent
It is guaranteed that the sum of $N$ over all testcases does not exceed $5 \cdot 10^5$.

\OutputFile

\noindent
For each testcase, output the answer in the following format:

\noindent
If a valid string exists, print \textbf{YES} on the first line. On the second line, print the valid string $S$ of length $N$ consisting of uppercase and lowercase English letters that satisfies all the given conditions. If there are multiple valid strings, any valid string will be accepted.

\noindent
If no such string exists, print \textbf{NO} on a single line.

\Examples

\begin{example}
\exmp{1}{4
2
11
4
1010
3
010
8
11101000
}{YES
xx
YES
axad
NO
YES
axxxadef
}
\end{example}

\pagebreak

\end{problem}